\uuid{YTqz}
\titre{Notion d'équation différentielle}
\auteur{bodin ; croizat ; sacre}
\organisation{exo7}
\licence{CC BY-NC-SA 4.0}
\datecreate{2022-03-08}
\niveau{L1}
\module{Analyse}
\chapitre{Equation différentielle}
\sousChapitre{Résolution d'équation différentielle du premier ordre}
\difficulte{2}
\theme{Notion d'équation différentielle}
\exo7qcmid{hilisit.equations-differentielles.017}

\contenu{
\begin{question}
\qclassification{225.01}
Soit l'équation différentielle $y'=2x(y+x)-1$. Quelles sont les affirmations vraies ?

\begin{answers}
 \good{$y= e^{x^2}-x$ est une solution.}
 \good{Cette équation différentielle n'a pas de solution constante.}
 \good{$y=-x$ est une solution.}
 \bad{$y= e^{x^2}-x+1$ est une solution.}
\end{answers}
\begin{explanations} Pour une fonction constante $y=C$, $y'=0$ et $2x(y+x)-1=2x(C+x)-1$, ce qui n'est pas la fonction nulle (c'est un polynôme du second degré), donc $y=C$ n'est pas solution. Pour $y=-x$, $2x(y+x)-1=-1=y'$, donc $y=-x$ est solution. Pour $y= e^{x^2}-x$, $2x(y+x)-1=2x e^{x^2}-1=y'$ donc $y= e^{x^2}-x$ est une solution. Pour $y= e^{x^2}-x+1$, $y'=2x e^{x^2}-1$ et  $2x(y+x)-1=2x e^{x^2}+2x-1$ donc $y= e^{x^2}-x+1$ n'est pas solution.
\end{explanations}
\end{question}
}
